\paragraph{The published prime-two Euler-system input.}
For $p=2$ we use the rational, away-from-two case of
\cite[Theorem 0.8, pp.~317--318]{kato1999}, with the stronger
hypothesis in its condition 0.6(ii), to justify the application of
\cite[Theorem 13.4(2)]{kato}. For our non-CM elliptic representation
the image of $G_{\Q^{\rm ab}}$ contains an open subgroup of
$\mathrm{SL}_2(\Z_2)$, as stated in
\cite[Example 0.6.2, pp.~317--318]{kato1999}, using Serre's open-image
theorem. A quadratic character $\chi$ is trivial on $G_{\Q^{\rm ab}}$.
Thus the same subgroup acts on $V\chi$ and contains a matrix
$\left(\begin{smallmatrix}1&b\\0&1\end{smallmatrix}\right)$ with
$b\ne0$. Its fixed space over $\Q_2$ has dimension one. This supplies
the unipotent in $G_{\Q^{\rm ab}}$, not merely in the cyclotomic
subgroup. The latter subgroup also contains this open
$\mathrm{SL}_2$ subgroup, so $V\chi$ remains absolutely irreducible
over $\Q(\mu_{2^\infty})$. Indeed, a line invariant under both upper
and lower nonidentity unipotents cannot exist, even over an algebraic
closure. No inference from irreducibility over $G_\Q$ alone is needed.

The retained classes (K8) satisfy the Euler-system norm relation
0.1.1 of \cite{kato1999}, with the same inverse-Frobenius polynomial
(K9). Enlarge the fixed inverted support to contain $2$, the bad
primes, and the conductor of $\chi$; primes excluded only by smoothing
belong to $N'$ in that source and need not be inverted. The preceding
nonvanishing argument supplies a nontorsion class on each sign
component, purity gives $w=-1$, and oddness gives condition
0.6(iv$_{\mathfrak p}$). At the height-one primes used here, $2$ is
invertible and the localized Iwasawa algebra is a DVR. Theorem 0.8
then bounds the coordinate ideal of this actual retained class. Its
strong residual and integral-cokernel assumptions concern primes
containing $2$, which are not used here.

We also account for the difference between the two published
cohomology conventions. Write $C_j$ for the $j_*$ Iwasawa complex of
\cite[Section 12.1]{kato}, and $C_S$ for full cochains on the open
scheme with the odd support inverted, as in \cite[Section 0.5]{kato1999}.
Work in one rational sign component and localize at $\mathfrak p$.
Their quotient is the direct sum of the odd local complexes (K16).
Each differential has nonzero determinant, so its degree-one
cohomology vanishes. Hence $H^1(C_j)\simeq H^1(C_S)$ and
$H^2(C_j)\hookrightarrow H^2(C_S)$. At an odd prime $q$, the
unramified local term has cohomological dimension one after inverting
$2$. The localization triangle identifies the degree-two cohomology
of (K16) with full local $H^2(\Q_q,T\chi)$ in Iwasawa cohomology.
The map from $H^2(C_S)$ is the usual local restriction. Exactness
therefore identifies $H^2(C_j)$ with the kernel of all those odd
local restrictions. Intersecting with the kernel of restriction at
$2$ gives
\[
 \ker\bigl(H^2(C_j)\to H^2_{\rm loc,2}\bigr)
 \simeq
 \ker\bigl(H^2(C_S)\to\!\bigoplus_{q\in S}H^2_{\rm loc,q}\bigr).
\]
This is precisely the comparison of the subscript-zero groups in
the two sources. The $H^1$ identification preserves the actual
Euler-system class and its coordinate ideal $J$. Thus Theorem 0.8
gives (K15) in the convention used here. The argument for (K16)
below is local and does not use (K15). We do not invoke the
odd-prime-only Perrin--Riou argument or an unchecked version of
Rubin's theorem at $2$.

